\chapter{Introduction}
\label{ch:introduction}

Particle physics is the study of the fundamental constituents of matter and the forces that govern their interactions. This chapter provides an overview of the field and the motivation for this research.

\section{Historical Background}

The journey to understand the nature of matter has been one of humanity's greatest intellectual pursuits. From the atomic theory of Democritus to the discovery of the electron by J.J. Thomson in 1897, our understanding has evolved dramatically \cite{thomson1897}.

\section{The Standard Model}

The Standard Model of particle physics is a quantum field theory that describes three of the four known fundamental forces (electromagnetic, weak, and strong interactions) and classifies all known elementary particles \cite{standardmodel}.

\subsection{Fermions}

Fermions are the matter particles, divided into:
\begin{itemize}
  \item Quarks: up, down, charm, strange, top, bottom
  \item Leptons: electron, muon, tau, and their corresponding neutrinos
\end{itemize}

\subsection{Bosons}

Bosons are the force carriers:
\begin{itemize}
  \item Photon ($\gamma$) - electromagnetic force
  \item W and Z bosons - weak force
  \item Gluons - strong force
  \item Higgs boson - mass generation
\end{itemize}

\section{Motivation for This Research}

Despite the remarkable success of the Standard Model, several open questions remain:
\begin{itemize}
  \item What is the nature of dark matter?
  \item Why is there more matter than antimatter in the universe?
  \item How can we unify all fundamental forces?
\end{itemize}

This thesis addresses [specific aspect] of these fundamental questions.

\section{Thesis Outline}

The remainder of this thesis is organized as follows:
\begin{itemize}
  \item Chapter \ref{ch:theory} presents the theoretical framework
  \item Chapter \ref{ch:experiment} describes the experimental setup
  \item Chapter \ref{ch:results} presents the results and analysis
  \item Chapter \ref{ch:conclusion} summarizes the conclusions and future outlook
\end{itemize}
